\documentclass[10pt,a4paper]{amsart}
\usepackage[margin=19mm]{geometry}
\usepackage{amsmath,amssymb,mathtools}
\usepackage[scr=rsfs,cal=euler]{mathalfa}
\usepackage{xfrac}
\usepackage[unicode,hidelinks,bookmarks=false]{hyperref}
\newcommand{\ie}{i.e.\ }
\newcommand{\cf}{cf.\ }
\newcommand{\resp}{respectively\ }
\newcommand{\Subsection}{\subsection}
\providecommand{\nobreakdash}{\nobreak\hbox{-}}
\setlength{\emergencystretch}{2em}
\multlinegap0pt
\usepackage{float}
\usepackage{amssymb}
\newtheorem{lemma}{Lemma}
\title{Exponential Decay \\for a Boundary-Controlled \\Nonlinear Parabolic Reactor Model}
\author{Yevgeniia Yevgenieva, Alexander Zuyev, Christophe Prieur, Peter Benner}
\date{Source: arXiv:2603.23151v3 (CC BY 4.0)}
\begin{document}
\maketitle

\noindent\emph{Source and licence.} Exponential Decay \\for a Boundary-Controlled \\Nonlinear Parabolic Reactor Model. Version arXiv:2603.23151v3 (CC BY 4.0), distributed by arXiv under the Creative Commons Attribution 4.0 International licence, http://creativecommons.org/licenses/by/4.0/, as recorded in the arXiv licence metadata for this exact version and shown on the abstract page https://arxiv.org/abs/2603.23151v3. Full licence notice: \texttt{licenses/arxiv-2603.23151-CC-BY-4.0.txt}.\par\smallskip

\noindent\textit{Editorial scope.} Counted unit: the article's lemma lem:dependelne\_on\_alpha (source0.tex lines 336--351). The article's complete original proof of it is delivered with it (source0.tex lines 353--447, one \texttt{proof} environment of 95 lines). No mathematics is written, repaired or re-derived: the statement and the proof are the article's own lines, in the article's own notation, and no retained line is substituted or re-flowed.  Retained uncounted as the source-local context the proof needs: lines 179--182; lines 307--309; lines 323--325; lines 332--334.  Nothing was omitted from any retained span, so the package carries no omission record.  The retained lines cite 1 works, whose entries are transcribed verbatim from the article's own bibliography source in the e-print and delivered with short editorial labels recorded in the item's reference mapping.  No retained line contains a figure environment, an image-inclusion command or drawing code, so no image is delivered and the package declares no asset dependency.  Editorial formatting only, in addition: an AMS \texttt{amsart} preamble that restates the article's own declarations and macros used by the retained lines, namely the environments \texttt{lemma} on separate counters, together with the standard packages those lines need.  The retained environments print in this document as Lemma 1 in order of appearance, whereas the article prints the same items with the numbering of its own full document.  The complete native source of the article as distributed in the arXiv e-print for this exact version is bundled unchanged as \texttt{sources/197056/source0.tex}.
\begin{equation}\label{opA}
    \mathcal{A}:{\xi}\in D(\mathcal{A}) \mapsto \mathcal{A}{\xi} =
            (D\xi'' -v\xi')\in X,
\end{equation}

\begin{equation}\label{eq:eigen_problem}
\mathcal{A}\xi=\lambda\xi,\qquad \lambda\in\mathbb{C},\ \xi\in D(\mathcal{A}).
\end{equation}

\begin{equation}\label{eq:operator_L}
\mathcal{L}y:=- D y'' = \theta y, \qquad y\in D(\mathcal{L}),
\end{equation}

\begin{equation}\label{eq:alpha*}
\alpha^*:=\tfrac{1}{2}+\tfrac{D}{vl+2D}\in\left(\tfrac{1}{2},1\right).   
\end{equation}

\begin{lemma}\label{lem:dependelne_on_alpha}
For the operator $\mathcal{A}$ defined by \eqref{opA}, the following statements hold:
\begin{itemize}
    \item[$(i)$] The principal eigenvalue $\lambda_0(\alpha)$ is increasing in $\alpha\in\mathbb R$.

    \item[$(ii)$] If $\alpha < \alpha^*$, then $\lambda_0(\alpha)<-\dfrac{v^2}{4D}$.

    \item[$(iii)$] If $\alpha = \alpha^*$, then $\lambda_0(\alpha)=-\dfrac{v^2}{4D}$.

    \item[$(iv)$] If $\alpha^*<\alpha<1$, then $-\dfrac{v^2}{4D}<\lambda_0(\alpha)<0$.
    
    \item[$(v)$] If $\alpha = 1$, then $\lambda_0(\alpha)=0$.
    
    \item[$(vi)$] If $\alpha > 1$, then $\lambda_0(\alpha)>0$.
\end{itemize}
\end{lemma}

\begin{proof}
Using the standard variational characterization of the principal eigenvalue of a self-adjoint Sturm--Liouville operator $\mathcal{L}$ by means of the Rayleigh quotient (see, e.g., \cite{Zet2005}), we get $\theta_0(\alpha)=\min_{0\neq y\in H^2(0,l)}\frac{\langle\mathcal{L}y,y\rangle}{\langle y,y\rangle}$, where $\langle\cdot,\cdot\rangle$ denotes the standard inner product in $L^2(0,l)$. This implies:
\begin{equation*}
\theta_0(\alpha)=\min_{\substack{y\in H^2(0,l)\\ y\neq 0}}\frac{D\int_0^l |y'|^2\,dx - D\delta |y(l)|^2 + D\gamma(\alpha)|y(0)|^2}
{\int_0^l |y|^2\,dx}.
\end{equation*}
We see that $\theta_0(\alpha)$ is increasing in $\gamma(\alpha)$, hence decreasing in $\alpha$.
Therefore, $\lambda_0(\alpha) = -\frac{v^2}{4D} - \theta_0(\alpha)$ is increasing in $\alpha$, which proves $(i)$.

Now, we focus on finding the eigenvalues of the operator $\mathcal{L}$, defined in \eqref{eq:operator_L}. Consider three cases.
\begin{enumerate}
    \item[1)] $\theta=0$, meaning that $\lambda=-\frac{v^2}{2D}$.
\end{enumerate}
Then, the solutions of \eqref{eq:operator_L} take the form
%\begin{equation*}
$y(x)=C_1x+C_2,$ 
%\; \forall 
for all $x \in (0,l)$,
%\qquad 
for constants $C_1,\,C_2\in\mathbb{R}$,
%\end{equation*}
and the boundary conditions imply:
%\begin{equation*}
$C_1-\gamma(\alpha)C_2=0$ and
$(1+\delta l)C_1+\delta C_2=0$.
%\end{equation*}
The corresponding determinant condition gives $\gamma(\alpha)=-\frac{\delta}{1+\delta l}$ which implies $\alpha=\alpha^*$. Hence, the eigenvalue $\lambda=-\frac{v^2}{2D}$ belongs to the spectrum of the operator $\mathcal{A}$ only in the case $\alpha=\alpha^*$, where $\alpha^*$ is defined in \eqref{eq:alpha*}.

\begin{enumerate}
    \item[2)] $\theta<0$. Denote $\theta=-Dq^2$ with some parameter $q>0$.
\end{enumerate}
 In this case, the solutions of the eigenvalue problem \eqref{eq:operator_L} have the following form: 
$
%\begin{equation*}
y(x)=C_1e^{qx}+C_2e^{-qx}$, for all 
$x \in (0,l)$ with $C_1,\,C_2\in\mathbb{R}$,
%\end{equation*}
and the boundary conditions yield:
\begin{equation}\label{eq:BC_nu<0}
(q-\gamma)C_1-(q+\gamma)C_2= 
(q+\delta)e^{2ql}C_1-(q-\delta)C_2=0.
\end{equation}
The corresponding determinant condition can be written as
\begin{equation}\label{eq:R(q)}
e^{2ql}=\frac{(q-\gamma)(q-\delta)}{(q+\gamma)(q+\delta)}=:R(q), \qquad q>0.
\end{equation}
At this stage, we assume that $q\neq -\gamma$ when $\gamma<0$, since this would imply
$\delta=-\gamma$, that is, $\alpha=1$. The case $\alpha=1$ will be treated separately below.

Now we analyze equation \eqref{eq:R(q)}. First, we note that $e^{2ql}>1$ since $q>0$. A standard analysis shows that $R(q)>1$ only if $\gamma<0$ $\left(\alpha>\frac{1}{2}\right)$ and $0<q<-\gamma$. Moreover, under the condition $\gamma<0$, the function $S(q):=R(q)-e^{2ql}$ satisfies $S(0)=0$ and $S(q)\to+\infty$ as $q\to-\gamma$. This means that a sufficient condition for the existence of roots $S(q)=0$ on the interval $q\in(0,-\gamma)$ is $S'(0)<0$, which leads to the condition $\alpha>\alpha^*$ with $\alpha^*$ from \eqref{eq:alpha*}.
%
%It is easy to show that in the case $\gamma\geqslant0$, $R(q)<1$ for any $q$, which means that equation \eqref{eq:R(q)} does not have any real roots. It leads to the conclusion that in the case $\gamma\geqslant0$ operator $\mathcal{L}$ does not have positive eigenvalues. This case corresponds to $\alpha\leqslant\frac{1}{2}$. This proves $(ii)$.
%
%In the case $\gamma<0$ $\left(\alpha>\frac{1}{2}\right)$, $R(q)$ can be grater than $1$, which means that $\lambda$ can be grater than $-\frac{v^2}{4D}$. Nevertheless, we note that in the case $\frac{1}{2}<\alpha< 1$, $\delta>-\gamma$ and therefore $R(q)>1$ only if $0<q<-\gamma$, which translates into
%\begin{equation*}
%\begin{aligned}
%\lambda_0(\alpha)&=-\frac{v^2}{4D}+Dq^2<-\frac{v^2}{4D}+D\,\gamma^2
%\\&=\frac{v^2}{D}\left(\Big(\alpha-\frac12\Big)^2-\frac{1}{4}\right)<0.
%\end{aligned}
%\end{equation*}
This proves $(ii)$ and $(iii)$.

\begin{enumerate}
    \item[3)] $\theta>0$. Denote $\theta=Ds^2$ with some parameter $s>0$.
\end{enumerate}
 In this case, the solutions of the eigenvalue problem \eqref{eq:eigen_problem} have the following form:
%\begin{equation*}
$y(x)=C_1\cos(sx)+C_2\sin(sx)$,
for all $x \in (0,l)$. 
%\end{equation*}
The determinant condition for the 
%corresponding 
boundary conditions lead to 
%the following equation for $q>0$:
% \begin{equation*}
% \begin{aligned}
% \gamma\,C_1-qC_2&=0,
% \\(\delta\cos(ql)-q\sin(ql))C_1+(\delta\sin(ql)+q\cos(ql))C_2&=0,
% \end{aligned}
% \end{equation*}
% which leads to the transcendental equation for $q$:
\begin{equation}\label{eq:for_s}
(s^2-\delta\gamma(\alpha))\tan(sl)=s(\gamma(\alpha)+\delta),\qquad s>0.
\end{equation}

\begin{enumerate}
    \item[4)] $\alpha=1$. In this case, $\gamma=-\delta$.
\end{enumerate}
The boundary conditions \eqref{eq:BC_nu<0} translate into:
$
(q+\delta)C_1=0$, $(q+\delta)e^{ql}C_1-(q-\delta)e^{-ql}C_2=0$,
which leads to the only positive eigenvalue $\theta=-D\delta^2=-\frac{v^2}{4D}$ for operator $\mathcal{L}$, which proves $(v)$.

Finally, statements $(iv)$ and $(vi)$ follow from the others.
\end{proof}

\begin{thebibliography}{99}
\bibitem[Zet200]{Zet2005} {A.~Zettl}, {\em Sturm-{L}iouville theory}, vol.~121 of Mathematical Surveys and Monographs, American Mathematical Society, Providence, RI, 2005.
\end{thebibliography}
\end{document}
